\begin{mistake}{2026-05-28}{04}

\begin{problemcontent}
已知
\[
f(x,y)=
\begin{cases}
xy\dfrac{x^2-y^2}{x^2+y^2},&(x,y)\ne(0,0),\\[6pt]
0,&(x,y)=(0,0),
\end{cases}
\]
则下列结论正确的是（ ）。
\[
\begin{array}{ll}
(A)\ f''_{xy}(0,0)=1, & (B)\ f''_{xy}(0,0)=0,\\[4pt]
(C)\ f''_{yx}(0,0)=1, & (D)\ f''_{yx}(0,0)=-1.
\end{array}
\]
\end{problemcontent}

\begin{solutioncontent}
采用记号 \(f''_{xy}=(f_x)_y\)，\(f''_{yx}=(f_y)_x\)。由于在坐标轴上
\[
f(x,0)=0,\qquad f(0,y)=0,
\]
需要从偏导数定义出发计算原点处的混合偏导数。

先求 \(f_x(0,y)\)。当 \(y\ne 0\) 时，
\[
\begin{aligned}
f_x(0,y)
&=\lim_{h\to 0}\frac{f(h,y)-f(0,y)}{h}\\
&=\lim_{h\to 0}\frac{hy\dfrac{h^2-y^2}{h^2+y^2}}{h}\\
&=\lim_{h\to 0}y\frac{h^2-y^2}{h^2+y^2}
=-y.
\end{aligned}
\]
当 \(y=0\) 时，\(f_x(0,0)=0\)，故 \(f_x(0,y)=-y\)。于是
\[
\begin{aligned}
f''_{xy}(0,0)
&=(f_x)_y(0,0)\\
&=\lim_{k\to 0}\frac{f_x(0,k)-f_x(0,0)}{k}\\
&=\lim_{k\to 0}\frac{-k-0}{k}=-1.
\end{aligned}
\]

再求 \(f_y(x,0)\)。当 \(x\ne 0\) 时，
\[
\begin{aligned}
f_y(x,0)
&=\lim_{k\to 0}\frac{f(x,k)-f(x,0)}{k}\\
&=\lim_{k\to 0}\frac{xk\dfrac{x^2-k^2}{x^2+k^2}}{k}\\
&=\lim_{k\to 0}x\frac{x^2-k^2}{x^2+k^2}
=x.
\end{aligned}
\]
当 \(x=0\) 时，\(f_y(0,0)=0\)，故 \(f_y(x,0)=x\)。于是
\[
\begin{aligned}
f''_{yx}(0,0)
&=(f_y)_x(0,0)\\
&=\lim_{h\to 0}\frac{f_y(h,0)-f_y(0,0)}{h}\\
&=\lim_{h\to 0}\frac{h-0}{h}=1.
\end{aligned}
\]

因此
\[
f''_{xy}(0,0)=-1,\qquad f''_{yx}(0,0)=1,
\]
正确选项为 \((C)\)。
\end{solutioncontent}

\mistakeknowledge{
\begin{itemize}
  \item \textbf{核心公式/定理}：混合偏导按定义计算，\(f''_{xy}=(f_x)_y\)，\(f''_{yx}=(f_y)_x\)；只有在相应二阶偏导连续等条件下，才能直接使用混合偏导相等定理。
  \item \textbf{易错点}：原点处函数为分段定义，不能直接对 \(xy\dfrac{x^2-y^2}{x^2+y^2}\) 求偏导后代入 \((0,0)\)，也不能默认 \(f''_{xy}(0,0)=f''_{yx}(0,0)\)。
  \item \textbf{关联知识}：求原点处混合偏导时，常先用定义求轴线上函数 \(f_x(0,y)\)、\(f_y(x,0)\)，再继续对另一变量求导。
\end{itemize}}

\end{mistake}
